\documentclass{article}
\usepackage[utf8]{inputenc}
\usepackage{amsmath, amssymb, amsthm}
\usepackage{bm}
\usepackage{hyperref}
\usepackage{geometry}
\geometry{a4paper, margin=1in}

\title{[minireport]BoundedDenoiser}
\author{jeon.isavelle@gmail.com}
\date{\today}

\begin{document}

\maketitle

\section{Introduction}
While a deep neural network $D_{\sigma}$ can act as an effective denoiser, an arbitrary architecture lacks the necessary regularity for iterative optimization.
As the noise level $\sigma$ decreases, the denoiser often exhibits significant instability, where artifacts are amplified over successive iterations.
In frameworks such as Half-Quadratic Splitting (HQS) or ADMM, a lack of structural constraints prevents the algorithm from reaching a stable fixed point, leading to degraded image quality.

Among representative formulations including such as bounded denoiser, nonexpensive denoiser, averaged operator, the methodology of bounded denoiser will be experimented and proved.
This choice is motivated by following advantages:

\begin{itemize}
    \item it enables the stable implementation of diminishing-$\sigma$ strategies.
    \item it provides fixed-point convergence guarantees even in nonconvex and ill-posed inverse problems, such as Poisson denoising and blind deconvolution, where the data-fidelity term may be non-Lipschitz or not strongly convex.
\end{itemize}

\subsection{Stable Implementation of Diminishing $\sigma$ Strategies}
Lipschitz-constrained denoisers make it possible to implement diminishing-$\sigma$ continuation strategies in a stable and controlled manner.
This is a key requirement for adaptive PnP-ADMM, particularly in nonconvex and ill-posed settings such as Poisson denoising and blind deconvolution, where the data-fidelity term may fail to satisfy standard assumptions such as Lipschitz continuity or strong convexity.




\section{Background}
To analyze and stabilize the iterative trajectory of Plug-and-Play (PnP) methods, several geometric and operator-theoretic quantities must be characterized at each iteration.
In particular, convergence behavior is closely related to the regularity of the denoiser, the structure of the PnP-ADMM update, and the local sensitivity of the denoiser as measured by its Jacobian.

\subsection{Theory of Lipschitz Convergence}
A standard route to provable convergence in PnP methods is to impose a Lipschitz constraint on the denoiser.
If the overall iteration becomes contractive, the Banach contraction principle guarantees convergence to a unique fixed point, typically with linear (geometric) convergence rate.

A common requirement is that the denoiser be nonexpansive, or preferably contractive:
\begin{equation}
    \|D(x) - D(y)\|_2 \le L \|x - y\|_2,
\end{equation}
where $0 \le L \le 1$.
When $L \le 1$, the operator is nonexpansive; when $L < 1$, it is contractive.

Such regularity is particularly important in iterative inverse problems, where repeated application of an unstable denoiser can amplify perturbations and prevent convergence.

\subsection{PnP-ADMM}
Consider the inverse problem
\begin{equation}
    \min_x f(x) + g(x),
\end{equation}
where $f(x)$ is a data-fidelity term and $g(x)$ is a regularization term.
In Plug-and-Play ADMM, the proximal operator associated with $g$ is replaced by a learned denoiser $D_\sigma$.
The resulting iteration is
\begin{align}
    x^{k+1} &= \arg\min_x \; f(x) + \frac{\rho}{2}\|x - v^k + u^k\|_2^2, \\
    v^{k+1} &= D_\sigma(x^{k+1} + u^k), \\
    u^{k+1} &= u^k + x^{k+1} - v^{k+1},
\end{align}
where $\rho > 0$ is the penalty parameter and $u^k$ is the scaled dual variable.

The convergence of this scheme depends critically on the regularity of $D_\sigma$.
If the denoiser is excessively sensitive, especially at small noise levels, the iteration may fail to approach a stable fixed point.

\subsection{Jacobian-Vector Product}
Let $f : \mathbb{R}^n \to \mathbb{R}^m$ be a differentiable function and let $x \in \mathbb{R}^n$.
The Jacobian matrix of $f$ at $x$ is defined by
\begin{equation}
    J_{ij}(x) = \frac{\partial f_i}{\partial x_j}(x).
\end{equation}

For an arbitrary vector $v \in \mathbb{R}^n$, the Jacobian-vector product (JVP) is defined as
\begin{equation}
    J(x)v
    = \left.\frac{d}{d\epsilon} f(x + \epsilon v)\right|
    = \lim_{\epsilon \to 0} \frac{f(x + \epsilon v) - f(x)}{\epsilon}.
\end{equation}

This quantity is exactly the directional derivative of $f$ at $x$ along direction $v$.
The goal is to improve the stability of the denoiser used inside PnP iterations by controlling its local sensitivity and approximate Lipschitz behavior.

\subsection{Jacobian Regularization}
One practical strategy is to penalize the magnitude of Jacobian-vector products during training.
Adding a regularization term of the form $\|Jv\|_2^2$ encourages the denoiser to behave more smoothly with respect to input perturbations, which can substantially improve optimization stability.

We regularize the Jacobian of the residual map
\begin{equation}
    R(x) = x - H(x),
\end{equation}
where $H$ denotes the denoiser or restoration operator.

A computationally efficient estimator can be constructed using Hutchinson's trick together with JVPs:
\begin{equation}
    \mathcal{L}_{\mathrm{jac}}
    = \mathbb{E}_{x,z}\left[\|J_R(x) z\|_2^2\right],
    \qquad z \sim \mathcal{N}(0, I).
\end{equation}

This penalty discourages large local amplification in random directions and therefore promotes smoother, more stable operator behavior.

\subsection{Real Spectral Normalization}
As a complementary strategy, Real Spectral Normalization (RealSN) aims to bound the approximate Lipschitz constant of the denoiser, preferably through its residual operator.
By controlling the operator norm, RealSN promotes nonexpansive or contraction-like behavior, which is desirable for stabilizing iterative PnP updates.

From the perspective of fixed-point theory, this form of spectral control is attractive because it directly targets the mechanism responsible for iterative instability: local expansion of perturbations under repeated denoising steps.
Therefore, Jacobian regularization and RealSN can be viewed as two complementary tools for improving convergence robustness in PnP reconstruction.

\subsection{Why Regularize the Residual Map?}
In practice, it is often more natural to control the residual map
\begin{equation}
    R(x) = x - H(x)
\end{equation}
rather than directly regularizing the denoiser $H$ itself.

The main reason is that the residual map isolates the \emph{correction} applied by the denoiser relative to the identity operator.

From a differential viewpoint, the Jacobian of the residual map is
\begin{equation}
    J_R(x) = I - J_H(x).
\end{equation}
Therefore, controlling $J_R(x)$ amounts to controlling how far the denoiser deviates locally from the identity mapping.
This is particularly relevant in PnP algorithms, where overly aggressive denoising may destroy data-consistent structure and destabilize the fixed-point iteration.

\subsection{Connection to PnP Stability}
By encouraging smoother residual corrections and bounding the effective Lipschitz behavior of the denoiser, Jacobian regularization and RealSN improve the stability of the composed iterative map.
This makes diminishing-$\sigma$ schedules more reliable and better suited to ill-posed inverse problems.


\section{Protocol}

\subsection{Regularization and Stabilization Settings}
We evaluate two complementary stabilization strategies for the denoiser used in PnP-ADMM: Jacobian regularization and Real Spectral Normalization (RealSN).

For Jacobian regularization, we consider the penalty
\begin{equation}
    \mathcal{L}_{\mathrm{jac}}
    = \mathbb{E}_{x,z}\left[\|J_R(x)z\|_2^2\right],
    \qquad z \sim \mathcal{N}(0, I),
\end{equation}
where $R(x) = x - H(x)$ denotes the residual map.
In practice, the following hyperparameter settings are used:
\begin{itemize}
    \item Jacobian regularization weight: $\beta \in \{0, 10^{-4}, 10^{-3}, 10^{-2}\}$,
    \item number of random probe vectors: \texttt{z\_samples} $=1$ per batch.
\end{itemize}
Using a single probe vector per batch provides a computationally efficient estimator while keeping the notebook-based experiments tractable.

For Real Spectral Normalization, we follow a two-phase implementation strategy.
In Phase 1, we use PyTorch's spectral normalization parametrization on convolutional layers for fast and robust experimentation.
In Phase 2, we implement a more paper-faithful RealSN variant, where each convolution is treated as a linear operator and its dominant singular value $\sigma(K_\ell)$ is estimated via power iteration and renormalization.

The main experimental knobs for RealSN:
\begin{itemize}
    \item spectral normalization: on/off,
    \item global scaling factor: $c \in \{1.0, 0.95, 0.9, 0.8\}$.
\end{itemize}


For analysis, we log the following quantities:
\begin{enumerate}
    \item Jacobian term: $\mathrm{mean}(\|J_R(x)z\|_2^2)$,
    \item per-layer spectral norm estimates: $\sigma(K_\ell)$,
    \item maximum and mean spectral norms across layers,
    \item the effective scaling budget $c_\ell$ applied to each layer.
\end{enumerate}

\subsection{Data and Split}
Experiments are conducted on a CT dataset.
The dataset is divided into training, validation, and test sets, where the validation split is used for hyperparameter selection and the test split is used only for final evaluation.
When synthetic inverse problems are considered, measurements are generated from clean images $x_0$ using a predefined forward operator and noise model, so that reconstruction quality can be evaluated against ground truth.

\subsection{Experimental Procedure}
The overall experimental pipeline is as follows:
\begin{enumerate}
    \item define the inverse problem and measurement model;
    \item load a pretrained denoiser;
    \item run the PnP-ADMM iteration under a selected continuation schedule;
    \item evaluate convergence behavior, reconstruction quality, and stability-related proxy metrics.
\end{enumerate}

We also study continuation schedules for the ADMM penalty parameter $\rho$.
When continuation is enabled, $\rho$ is updated according to
\begin{equation}
    \rho_{k+1} = \gamma \rho_k.
\end{equation}
For ViT- or DiT-like denoisers, we use $\gamma = 1.01$ as the default setting, since smaller increments often reduce oscillatory behavior.
For CNN- or U-Net-based denoisers, values in the range
\begin{equation}
    \gamma \in [1.05, 1.1]
\end{equation}
often work well in practice.
A small grid search over $\gamma$ is performed in \texttt{03\_ablation\_grid.ipynb}.

\subsection{Evaluation Axes and Metrics}
We evaluate the method along three primary axes:
\begin{itemize}
    \item spectral normalization: off vs.\ on, including the scaling sweep over $c$,
    \item Jacobian regularization: sweep over $\beta$,
    \item continuation schedule: sweep over $\gamma$.
\end{itemize}

The evaluation metrics are grouped into three categories.

\paragraph{Convergence.}
We measure:
\begin{itemize}
    \item final primal residual $r^k$,
    \item number of iterations required to reach a prescribed tolerance,
    \item divergence or oscillation rate across runs.
\end{itemize}

\paragraph{Reconstruction quality.}
For synthetic experiments where the clean target $x_0$ is available, we report:
\begin{itemize}
    \item PSNR,
    \item SSIM.
\end{itemize}

\paragraph{Bound-related proxy metrics.}
To analyze the stabilization effect of the proposed regularizers, we log:
\begin{itemize}
    \item maximum and mean layer spectral norms,
    \item magnitude of the Jacobian regularization term.
\end{itemize}

This design yields a clear trade-off analysis:
spectral normalization improves global stability, Jacobian regularization suppresses local jitter and oscillatory updates, and their combination is expected to improve both the reliability and speed of convergence toward a stable solution $x^\ast$.

\subsection{Expected Outcome}
We expect the proposed stabilization strategies to reduce oscillatory behavior in PnP-ADMM and improve convergence toward a fixed point.
In particular, the discrepancy between the primal variable and the denoised auxiliary variable is expected to decrease over iterations, i.e.,
\begin{equation}
    \|x^k - z^k\|_2 \rightarrow 0.
\end{equation}
At the same time, we expect improved reconstruction quality and more stable behavior under diminishing-noise or continuation-based schedules.

\section{Limitations}
Despite the expected stability benefits, several limitations remain.

\begin{itemize}
    \item stronger regularization may improve stability at the cost of reconstruction sharpness.
    If $\gamma$ is too large, the algorithm may terminate or become numerically unstable before the denoiser fully exhibits the intended bounded or contractive behavior.
    In such cases, the practical benefit of denoiser regularization may be weakened.
    \item proxy quantities such as Jacobian norms or layerwise spectral norms do not by themselves constitute a full convergence proof for the entire nonlinear PnP iteration.
\end{itemize}




\end{document}